% Source: Weinberg GR, physical PDF pp. 248--252
%         (printed pp. 225--229).
% Coverage: Section 9.4, Multipole Fields; equations (9.4.1)--(9.4.36).
% Figures/tables/footnotes: no figures or tables.
% Status: source-reviewed and compile-clean.
% Uncertainties: Weinberg's stacked perturbative-order numerals are rendered
%                as parenthesized superscripts.  The source prints d^4x' in
%                (9.4.29) and the first line of (9.4.33); these are corrected
%                to the spatial volume element d^3x', as required by the
%                immediately following angular/radial reduction.

\subsection{Multipole Fields}\label{sec:9.4}

As a first example, let us calculate the gravitational field far away from
an arbitrary finite distribution of energy and momentum. Let
\(T^{\mu\nu}(\mathbf x,t)\) vanish for \(r>R\), where
\(r=|\mathbf x|\). We may then expand the denominators
\(|\mathbf x-\mathbf x'|\) in Eqs.~\eqref{eq:9.1.59},
\eqref{eq:9.1.62}, and \eqref{eq:9.1.65} in inverse powers of \(r/R\):
\begin{equation}
  \frac1{|\mathbf x-\mathbf x'|}
  \longrightarrow
  \frac1r+\frac{\mathbf x\cdot\mathbf x'}{r^3}+\cdots.
  \tag{9.4.1}\label{eq:9.4.1}
\end{equation}
We find
\begin{equation}
  \phi
  \longrightarrow
  -\frac{G M^{(0)}}r
  -\frac{G\,\mathbf x\cdot\mathbf D^{(0)}}{r^3}
  +\order(r^{-3}),
  \tag{9.4.2}\label{eq:9.4.2}
\end{equation}
\begin{equation}
  \zeta_i
  \longrightarrow
  -\frac{4G P_i^{(1)}}r
  -\frac{2Gx^jJ_{ji}^{(1)}}{r^3}
  +\order(r^{-3}),
  \tag{9.4.3}\label{eq:9.4.3}
\end{equation}
\begin{equation}
  \psi
  \longrightarrow
  -\frac{G M^{(2)}}r
  -\frac{G\,\mathbf x\cdot\mathbf D^{(2)}}{r^3}
  +\order(r^{-3}),
  \tag{9.4.4}\label{eq:9.4.4}
\end{equation}
where
\begin{equation}
  M^{(0)}\equiv\int T^{00(0)}\,\dd^3x,
  \tag{9.4.5}\label{eq:9.4.5}
\end{equation}
\begin{equation}
  \mathbf D^{(0)}
  \equiv\int\mathbf x\,T^{00(0)}\,\dd^3x,
  \tag{9.4.6}\label{eq:9.4.6}
\end{equation}
\begin{equation}
  P^{i(1)}\equiv\int T^{i0(1)}\,\dd^3x,
  \tag{9.4.7}\label{eq:9.4.7}
\end{equation}
\begin{equation}
  J_{ij}^{(1)}
  \equiv2\int x^iT^{j0(1)}\,\dd^3x,
  \tag{9.4.8}\label{eq:9.4.8}
\end{equation}
\begin{equation}
  M^{(2)}
  \equiv\int\left(T^{00(2)}+T^{ii(2)}\right)\dd^3x,
  \tag{9.4.9}\label{eq:9.4.9}
\end{equation}
\begin{equation}
  \mathbf D^{(2)}
  \equiv
  \int\mathbf x
  \left(
    T^{00(2)}+T^{ii(2)}
    +\frac1{4\pi G}\partial_t^2\phi
  \right)\dd^3x.
  \tag{9.4.10}\label{eq:9.4.10}
\end{equation}
(The term \(\partial_t^2\phi\) does not contribute to \(M^{(2)}\),
because it equals \(-\boldsymbol\nabla\cdot(\partial_t\boldsymbol\zeta)/4\)
and hence vanishes upon integration.)

The field \(\psi\) has physical effects only through its presence in the
expansion
\[
  g_{00}=-1-2\phi-2\psi-2\phi^2+\order(\bar v^6).
\]
We may take account of it simply by replacing \(\phi\) everywhere with
\(\phi+\psi\). Within the accuracy of the post-Newtonian approximation,
\begin{equation}
  g_{00}
  =
  -1-2(\phi+\psi)-2(\phi+\psi)^2+\order(\bar v^6).
  \tag{9.4.11}\label{eq:9.4.11}
\end{equation}
Equations~\eqref{eq:9.4.2} and \eqref{eq:9.4.4} give the physically
significant field as
\begin{equation}
  \phi+\psi
  \longrightarrow
  -\frac{GM}{r}
  -\frac{G\,\mathbf x\cdot\mathbf D}{r^3}
  +\order(r^{-3}),
  \tag{9.4.12}\label{eq:9.4.12}
\end{equation}
where
\begin{equation}
  M\equiv M^{(0)}+M^{(2)},
  \qquad
  \mathbf D\equiv\mathbf D^{(0)}+\mathbf D^{(2)}.
  \tag{9.4.13}\label{eq:9.4.13}
\end{equation}
The quantity \(\mathbf D\) does not represent an effect of physical
importance, but rather just a displacement of the whole field, for
Eq.~\eqref{eq:9.4.12} may be written
\begin{equation}
  \phi+\psi
  \longrightarrow
  -\frac{GM}{|\mathbf x-\mathbf D/M|}
  +\order(r^{-3}).
  \tag{9.4.14}\label{eq:9.4.14}
\end{equation}
We could have avoided the \(\mathbf D\) term altogether by defining our
coordinate system with its origin at the center of energy. On the other
hand, the \(1/r\) and \(1/r^2\) terms in Eq.~\eqref{eq:9.4.3} are true
physical effects of great interest.

Useful properties of the moments of \(T^{\mu\nu}\) follow from energy and
momentum conservation. Equation~\eqref{eq:9.3.2} gives, in general,
\begin{equation}
  \frac{\dd M^{(0)}}{\dd t}=0,
  \tag{9.4.15}\label{eq:9.4.15}
\end{equation}
\begin{equation}
  \frac{\dd\mathbf D^{(0)}}{\dd t}=\mathbf P^{(1)}.
  \tag{9.4.16}\label{eq:9.4.16}
\end{equation}
If the energy-momentum tensor is time independent, then
\[
  \partial_iT^{i0(1)}=0,
\]
and integration by parts gives
\begin{equation}
  0
  =
  \int x^i\partial_jT^{j0(1)}\,\dd^3x
  =
  -P^{i(1)},
  \tag{9.4.17}\label{eq:9.4.17}
\end{equation}
\begin{equation}
  0
  =
  2\int x^ix^j\partial_kT^{k0(1)}\,\dd^3x
  =
  -J_{ij}^{(1)}-J_{ji}^{(1)}.
  \tag{9.4.18}\label{eq:9.4.18}
\end{equation}
The vanishing of \(\mathbf P^{(1)}\) for a static system is hardly
surprising. The antisymmetry of \(J_{ij}^{(1)}\) is less obvious and allows
us to write
\begin{equation}
  J_{ij}^{(1)}=\varepsilon_{ijk}J_k^{(1)},
  \tag{9.4.19}\label{eq:9.4.19}
\end{equation}
where \(\mathbf J^{(1)}\) is the angular-momentum vector
\begin{equation}
  J_k^{(1)}
  =
  \frac12\varepsilon_{ijk}J_{ij}^{(1)}
  =
  \int\varepsilon_{ijk}x^iT^{j0(1)}\,\dd^3x.
  \tag{9.4.20}\label{eq:9.4.20}
\end{equation}
Here \(\varepsilon^{123}=+1\) is the spatial Levi--Civita tensor. Using
Eqs.~\eqref{eq:9.4.17} and \eqref{eq:9.4.19} in
Eq.~\eqref{eq:9.4.3} gives
\begin{equation}
  \boldsymbol\zeta
  \longrightarrow
  \frac{2G}{r^3}(\mathbf x\times\mathbf J)
  +\order(r^{-3}).
  \tag{9.4.21}\label{eq:9.4.21}
\end{equation}

The results \eqref{eq:9.4.14} and \eqref{eq:9.4.21} hold generally only
far from the gravitating mass. However, they also hold right down to the
surface of a spherical distribution of energy and momentum. First suppose
that \(T^{\mu\nu}(\mathbf x,t)\) depends on position only through
\(r=|\mathbf x|\). For \(r>r'\), the angular average is
\[
  \frac1{4\pi}\int\frac{\dd\Omega'}{|\mathbf x-\mathbf x'|}
  =
  \frac12\int_0^\pi
  \frac{\sin\theta\,\dd\theta}
  {(r^2-2rr'\cos\theta+r'^2)^{1/2}}
  =
  \frac1r.
\]
Hence everywhere outside the sphere,
\begin{equation}
  \phi=-\frac{GM^{(0)}}r,
  \tag{9.4.22}\label{eq:9.4.22}
\end{equation}
\begin{equation}
  \boldsymbol\zeta=-\frac{4G\mathbf P^{(1)}}r,
  \tag{9.4.23}\label{eq:9.4.23}
\end{equation}
\begin{equation}
  \psi=-\frac{GM^{(2)}}r.
  \tag{9.4.24}\label{eq:9.4.24}
\end{equation}
If the sphere is at rest, \(\mathbf P^{(1)}\) vanishes, and
Eqs.~\eqref{eq:9.1.57}, \eqref{eq:9.1.60}, \eqref{eq:9.1.61},
\eqref{eq:9.1.63}, and \eqref{eq:9.4.13} give
\begin{equation}
  g_{00}\simeq
  -1+\frac{2MG}{r}-\frac{2M^2G^2}{r^2},
  \tag{9.4.25}\label{eq:9.4.25}
\end{equation}
\begin{equation}
  g_{i0}\simeq0,
  \tag{9.4.26}\label{eq:9.4.26}
\end{equation}
\begin{equation}
  g_{ij}\simeq
  \delta_{ij}+\frac{2MG}{r}\delta_{ij}.
  \tag{9.4.27}\label{eq:9.4.27}
\end{equation}
{\setlength{\emergencystretch}{1em}%
This agrees with the exact Schwarzschild solution, given in harmonic
coordinates by Eq.~\eqref{eq:8.2.15}:\par}
\[
  g_{00}=-\frac{1-MG/r}{1+MG/r},
  \qquad
  g_{i0}=0,
\]
\[
  g_{ij}
  =
  \left(1+\frac{MG}{r}\right)^2\delta_{ij}
  {}+
  \left(\frac{MG}{r}\right)^2
  \frac{1+MG/r}{1-MG/r}\frac{x^ix^j}{r^2}.
\]
There is an important difference in the two derivations: the exact
Schwarzschild solution was derived for a static, spherically symmetric
system, while the post-Newtonian solution is valid for a system that can
vary over times of order \(\bar r/\bar v\). It will be shown in
Section~\ref{sec:11.7} that the Schwarzschild solution is actually valid
outside any spherically symmetric system, static or not.

Now consider a system that is at rest and spherically symmetric, but rotates
with angular frequency \(\boldsymbol\omega(r)\). Its momentum density is
\begin{equation}
  T^{i0(1)}(\mathbf x',t)
  =
  T^{00(0)}(r')
  [\boldsymbol\omega(r')\times\mathbf x']_i.
  \tag{9.4.28}\label{eq:9.4.28}
\end{equation}
Equation~\eqref{eq:9.1.62} gives
\begin{equation}
  \boldsymbol\zeta(\mathbf x)
  =
  -4G\int
  \frac{\boldsymbol\omega(r')\times\mathbf x'}
  {|\mathbf x-\mathbf x'|}
  T^{00(0)}(r')\,\dd^3x'.
  \tag{9.4.29}\label{eq:9.4.29}
\end{equation}
The solid-angle integral is
\begin{equation}
  \int\frac{\dd\Omega'\,\mathbf x'}{|\mathbf x-\mathbf x'|}
  =
  \frac{4\pi r'^2}{3r^3}\,\mathbf x,
  \qquad r'<r,
  \tag{9.4.30}\label{eq:9.4.30}
\end{equation}
\begin{equation}
  \int\frac{\dd\Omega'\,\mathbf x'}{|\mathbf x-\mathbf x'|}
  =
  \frac{4\pi}{3r'}\,\mathbf x,
  \qquad r'>r.
  \tag{9.4.31}\label{eq:9.4.31}
\end{equation}
Thus the field outside the sphere is
\begin{equation}
  \boldsymbol\zeta(\mathbf x)
  =
  \frac{16\pi G}{3r^3}
  \left[
    \mathbf x\times
    \int\boldsymbol\omega(r')T^{00(0)}(r')r'^4\,\dd r'
  \right].
  \tag{9.4.32}\label{eq:9.4.32}
\end{equation}
The integral may be expressed in terms of the angular momentum:
\begin{equation}
  \begin{split}
  \mathbf J
  &=
  \int
  \mathbf x'\times
  [\boldsymbol\omega(r')\times\mathbf x']
  T^{00(0)}(r')\,\dd^3x'\\
  &=
  \int
  \left[
    r'^2\boldsymbol\omega(r')
    -\mathbf x'(\mathbf x'\cdot\boldsymbol\omega(r'))
  \right]T^{00(0)}(r')\,\dd^3x'\\
  &=
  \frac{8\pi}{3}
  \int\boldsymbol\omega(r')T^{00(0)}(r')r'^4\,\dd r'.
  \end{split}
  \tag{9.4.33}\label{eq:9.4.33}
\end{equation}
Consequently, everywhere outside the sphere,
\begin{equation}
  \boldsymbol\zeta(\mathbf x)
  =
  \frac{2G}{r^3}(\mathbf x\times\mathbf J),
  \tag{9.4.34}\label{eq:9.4.34}
\end{equation}
in agreement with Eq.~\eqref{eq:9.4.21}. Inside a hollow spinning sphere,
Eqs.~\eqref{eq:9.4.29} and \eqref{eq:9.4.31} give
\begin{equation}
  \boldsymbol\zeta(\mathbf x)=\mathbf x\times\boldsymbol\Omega,
  \tag{9.4.35}\label{eq:9.4.35}
\end{equation}
where
\begin{equation}
  \boldsymbol\Omega
  =
  \frac{16\pi G}{3}
  \int\boldsymbol\omega(r')T^{00(0)}(r')r'\,\dd r'.
  \tag{9.4.36}\label{eq:9.4.36}
\end{equation}
The implications of this result for Mach's principle are discussed in
Section~\ref{sec:9.7}.
